% !TEX program = lualatex
% ---------------------------------------------------------------------------------------------
%  Deformalization 2: local nilpotence over a cyclotomic tower   (riemann-rev)
%
%  Build:  python3 -B scripts/build_digest.py local-nilpotence
%
%  One part of a unit of the Lean archive behind L. Alpöge's announced proof of RH, rewritten as ordinary
%  mathematics. The prose is hand-written. Lean names, line ranges and counts come from data.tex; the quoted
%  Lean statements are copied from the archive by the script. The note was written from the reading project's
%  structured proofs; the Lean files of the unit were not opened (§ "Notes on the formal text").
%
%  Rules of the text: standard terminology; every symbol is explained where it first appears and listed again
%  in the notation at the end; no letter has two meanings. Letters follow the first note of the series: the
%  cyclotomic character is kappa (the archive's cycZ), U_n are the principal units; the archive's closed normal
%  subgroup U is H here, its levU l U n is H_n and its cycU l U is H_infinity.
%
%  Layout: the reading edition of the guide (11 pt Libertinus on a 140 mm measure); preamble as in the first note.
% ---------------------------------------------------------------------------------------------
\documentclass[11pt,a4paper]{article}
\ifdefined\pdfvariable\pdfvariable suppressoptionalinfo \numexpr 1+2+4+8\relax\fi
\usepackage[a4paper,hmargin=35mm,top=22mm,bottom=27mm,footskip=12mm,nohead]{geometry}

\usepackage{fontspec}
\directlua{
  luaotfload.add_fallback("gfallback", {"STIXTwoMath-Regular.otf:mode=node;", "DejaVuSans.ttf:mode=node;"})
  luaotfload.add_fallback("gmonofallback", {"DejaVuSansMono.ttf:mode=node;", "STIXTwoMath-Regular.otf:mode=node;"})
}
\setmainfont{LibertinusSerif}[Extension=.otf,UprightFont=*-Regular,ItalicFont=*-Italic,
  BoldFont=*-Semibold,BoldItalicFont=*-SemiboldItalic,RawFeature={fallback=gfallback}]
\setsansfont{Inter}[Extension=.otf,UprightFont=*-Regular,ItalicFont=*-Italic,
  BoldFont=*-SemiBold,BoldItalicFont=*-SemiBoldItalic,RawFeature={fallback=gfallback}]
\setmonofont{JetBrainsMono}[Extension=.otf,UprightFont=*-Regular,ItalicFont=*-Italic,
  BoldFont=*-SemiBold,BoldItalicFont=*-SemiBoldItalic,Scale=MatchLowercase,
  RawFeature={-calt;-liga;fallback=gmonofallback}]
\usepackage{amsmath}
\usepackage{amsthm}
\usepackage[math-style=TeX]{unicode-math}
\setmathfont{LibertinusMath-Regular.otf}

\usepackage{xcolor}
\usepackage{tikz}
\usetikzlibrary{calc,arrows.meta,positioning,backgrounds}
\usepackage[protrusion=true,expansion=true]{microtype}
\usepackage{enumitem}
\usepackage{needspace}
\usepackage{booktabs}
\usepackage{array}
\usepackage{longtable}
\usepackage{fancyhdr}
\usepackage{lastpage}
\usepackage{multicol}
\usepackage[hidelinks,pdfusetitle]{hyperref}

\definecolor{ink}{HTML}{0F172A}
\definecolor{inksoft}{HTML}{475569}
\definecolor{inkfaint}{HTML}{94A3B8}
\definecolor{rulegrey}{HTML}{CBD5E1}
\definecolor{arith}{HTML}{C2410C}
\definecolor{est}{HTML}{7C3AED}
\definecolor{ana}{HTML}{0891B2}
\definecolor{ok}{HTML}{047857}
\definecolor{leancol}{HTML}{3730A3}

\title{Local nilpotence over a cyclotomic tower}
\author{riemann-rev}
\date{}

% ------------------------------------------------------------------------------------- macros
\newcommand{\leanfont}{\ttfamily\fontsize{8.6}{10.5}\selectfont\color{leancol}}
\newcommand{\lean}[1]{\csname LN@#1\endcsname}
\newcommand{\leanquote}[1]{\csname LQ@#1\endcsname}
% generated by scripts/build_digest.py -- do not edit
\expandafter\newcommand\csname LN@XDual\endcsname{\href{https://github.com/Deicyde/riemann-hypothesis/blob/40e56bab026552e089d7afcfeff08d2defb9c005/src/Development/AutomorphicArithmetic/AutomorphicVarieties/GaloisLibrary/KummerTheory/Mk0IdMapCongr.lean\#L90-L99}{\leanfont Kum.XDual}}
\expandafter\newcommand\csname LN@Dinf\endcsname{\href{https://github.com/Deicyde/riemann-hypothesis/blob/40e56bab026552e089d7afcfeff08d2defb9c005/src/Development/AutomorphicArithmetic/AutomorphicVarieties/GaloisLibrary/KummerTheory/Mk0IdMapCongr.lean\#L195-L195}{\leanfont Kum.Dinf}}
\expandafter\newcommand\csname LN@Xinf\endcsname{\href{https://github.com/Deicyde/riemann-hypothesis/blob/40e56bab026552e089d7afcfeff08d2defb9c005/src/Development/AutomorphicArithmetic/AutomorphicVarieties/GaloisLibrary/KummerTheory/Mk0IdMapCongr.lean\#L197-L197}{\leanfont Kum.Xinf}}
\expandafter\newcommand\csname LN@funRep\endcsname{\href{https://github.com/Deicyde/riemann-hypothesis/blob/40e56bab026552e089d7afcfeff08d2defb9c005/src/Development/AutomorphicArithmetic/AutomorphicVarieties/GaloisLibrary/KummerTheory/Mk0IdMapCongr.lean\#L121-L124}{\leanfont Kum.funRep}}
\expandafter\newcommand\csname LN@conjCond\endcsname{\href{https://github.com/Deicyde/riemann-hypothesis/blob/40e56bab026552e089d7afcfeff08d2defb9c005/src/Development/AutomorphicArithmetic/AutomorphicVarieties/GaloisLibrary/KummerTheory/Mk0IdMapCongr.lean\#L103-L107}{\leanfont Kum.conjCond}}
\expandafter\newcommand\csname LN@EqHom\endcsname{\href{https://github.com/Deicyde/riemann-hypothesis/blob/40e56bab026552e089d7afcfeff08d2defb9c005/src/Development/AutomorphicArithmetic/AutomorphicVarieties/GaloisLibrary/KummerTheory/Mk0IdMapCongr.lean\#L109-L110}{\leanfont Kum.EqHom}}
\expandafter\newcommand\csname LN@Un\endcsname{\href{https://github.com/Deicyde/riemann-hypothesis/blob/40e56bab026552e089d7afcfeff08d2defb9c005/src/Development/AutomorphicArithmetic/AutomorphicVarieties/GaloisLibrary/KummerTheory/Mk0IdMapCongr.lean\#L186-L186}{\leanfont Kum.Un}}
\expandafter\newcommand\csname LN@funRepMem\endcsname{\href{https://github.com/Deicyde/riemann-hypothesis/blob/40e56bab026552e089d7afcfeff08d2defb9c005/src/Development/AutomorphicArithmetic/AutomorphicVarieties/GaloisLibrary/KummerTheory/FunRepMem.lean\#L5-L46}{\leanfont Kum.funRepMem}}
\expandafter\newcommand\csname LN@memPrinc\endcsname{\href{https://github.com/Deicyde/riemann-hypothesis/blob/40e56bab026552e089d7afcfeff08d2defb9c005/src/Development/AutomorphicArithmetic/AutomorphicVarieties/GaloisLibrary/KummerTheory/DinfNil.lean\#L7-L12}{\leanfont mem\_princ\_iff\_pow\_dvd\_sub\_one}}
\expandafter\newcommand\csname LN@powSucc\endcsname{\href{https://github.com/Deicyde/riemann-hypothesis/blob/40e56bab026552e089d7afcfeff08d2defb9c005/src/Development/AutomorphicArithmetic/AutomorphicVarieties/GaloisLibrary/KummerTheory/DinfNil.lean\#L14-L29}{\leanfont pow\_succ\_dvd\_pow\_prime\_sub\_one}}
\expandafter\newcommand\csname LN@qzTors\endcsname{\href{https://github.com/Deicyde/riemann-hypothesis/blob/40e56bab026552e089d7afcfeff08d2defb9c005/src/Development/AutomorphicArithmetic/AutomorphicVarieties/GaloisLibrary/KummerTheory/DinfNil.lean\#L31-L43}{\leanfont exists\_prime\_pow\_smul\_eq\_zero\_qz}}
\expandafter\newcommand\csname LN@princNhds\endcsname{\href{https://github.com/Deicyde/riemann-hypothesis/blob/40e56bab026552e089d7afcfeff08d2defb9c005/src/Development/AutomorphicArithmetic/AutomorphicVarieties/GaloisLibrary/KummerTheory/DinfNil.lean\#L45-L60}{\leanfont exists\_princ\_subset\_of\_mem\_nhds\_one}}
\expandafter\newcommand\csname LN@dinfNil\endcsname{\href{https://github.com/Deicyde/riemann-hypothesis/blob/40e56bab026552e089d7afcfeff08d2defb9c005/src/Development/AutomorphicArithmetic/AutomorphicVarieties/GaloisLibrary/KummerTheory/DinfNil.lean\#L62-L235}{\leanfont Kum.dinfNil}}
\expandafter\newcommand\csname LN@G\endcsname{\href{https://github.com/Deicyde/riemann-hypothesis/blob/40e56bab026552e089d7afcfeff08d2defb9c005/src/Formalism/ConjugateSelfDual/Statement.lean\#L6-L6}{\leanfont G}}
\expandafter\newcommand\csname LN@QZ\endcsname{\href{https://github.com/Deicyde/riemann-hypothesis/blob/40e56bab026552e089d7afcfeff08d2defb9c005/src/Formalism/GaloisCohomology/Statement.lean\#L15-L15}{\leanfont QZ}}
\expandafter\newcommand\csname LN@cycZ\endcsname{\href{https://github.com/Deicyde/riemann-hypothesis/blob/40e56bab026552e089d7afcfeff08d2defb9c005/src/ReciprocityTheory/GaloisValuations/Global.lean\#L90-L92}{\leanfont Gal.cycZ}}
\expandafter\newcommand\csname LN@Dec\endcsname{\href{https://github.com/Deicyde/riemann-hypothesis/blob/40e56bab026552e089d7afcfeff08d2defb9c005/src/ReciprocityTheory/GaloisValuations/Global.lean\#L15-L15}{\leanfont Gal.Dec}}
\expandafter\newcommand\csname LN@Ine\endcsname{\href{https://github.com/Deicyde/riemann-hypothesis/blob/40e56bab026552e089d7afcfeff08d2defb9c005/src/ReciprocityTheory/GaloisValuations/Global.lean\#L45-L64}{\leanfont Gal.Ine}}
\expandafter\newcommand\csname LN@princ\endcsname{\href{https://github.com/Deicyde/riemann-hypothesis/blob/40e56bab026552e089d7afcfeff08d2defb9c005/src/Development/AutomorphicArithmetic/CyclotomicBounds/ClassNumberBounds/TotallyRealRoots/LeNuIffPowMemPrinc.lean\#L16-L16}{\leanfont princ}}
\expandafter\newcommand\csname LN@levU\endcsname{\href{https://github.com/Deicyde/riemann-hypothesis/blob/40e56bab026552e089d7afcfeff08d2defb9c005/src/Development/AutomorphicArithmetic/AutomorphicVarieties/GaloisLibrary/SelmerTransport/ConjStableShaFam.lean\#L187-L187}{\leanfont Base.levU}}
\expandafter\newcommand\csname LN@cycU\endcsname{\href{https://github.com/Deicyde/riemann-hypothesis/blob/40e56bab026552e089d7afcfeff08d2defb9c005/src/Development/AutomorphicArithmetic/AutomorphicVarieties/GaloisLibrary/SelmerTransport/ConjStableShaFam.lean\#L189-L189}{\leanfont Base.cycU}}
\expandafter\newcommand\csname LN@conjHom\endcsname{\href{https://github.com/Deicyde/riemann-hypothesis/blob/40e56bab026552e089d7afcfeff08d2defb9c005/src/Development/AutomorphicArithmetic/AutomorphicVarieties/GaloisLibrary/SelmerTransport/RestrictionMemInvariantH1.lean\#L115-L118}{\leanfont Base.conjHom}}
\expandafter\newcommand\csname LN@conjDec\endcsname{\href{https://github.com/Deicyde/riemann-hypothesis/blob/40e56bab026552e089d7afcfeff08d2defb9c005/src/Development/AutomorphicArithmetic/AutomorphicVarieties/PartsConj.lean\#L69-L75}{\leanfont conjGal\_mem\_Dec}}
\expandafter\newcommand\csname LN@conjIne\endcsname{\href{https://github.com/Deicyde/riemann-hypothesis/blob/40e56bab026552e089d7afcfeff08d2defb9c005/src/Development/AutomorphicArithmetic/AutomorphicVarieties/PartsConj.lean\#L77-L81}{\leanfont conjGal\_mem\_Ine}}
\expandafter\newcommand\csname LN@xUnip\endcsname{\href{https://github.com/Deicyde/riemann-hypothesis/blob/40e56bab026552e089d7afcfeff08d2defb9c005/src/Development/AutomorphicArithmetic/AutomorphicVarieties/AdditiveValuationNorms/XUnip.lean\#L9-L27}{\leanfont PropA.xUnip}}
\expandafter\newcommand\csname LN@dualUnip\endcsname{\href{https://github.com/Deicyde/riemann-hypothesis/blob/40e56bab026552e089d7afcfeff08d2defb9c005/src/Development/AutomorphicArithmetic/AutomorphicVarieties/AdditiveValuationNorms/DualUnip.lean\#L9-L34}{\leanfont PropA.dualUnip}}
\expandafter\newcommand\csname LN@shaDGen\endcsname{\href{https://github.com/Deicyde/riemann-hypothesis/blob/40e56bab026552e089d7afcfeff08d2defb9c005/src/Development/AutomorphicArithmetic/AutomorphicVarieties/NumberFieldTowers/ShaDGen.lean\#L9-L59}{\leanfont Tower.shaDGen}}
\newcommand{\Pin}{40e56bab}
\newcommand{\PinUrl}{https://github.com/Deicyde/riemann-hypothesis/tree/40e56bab026552e089d7afcfeff08d2defb9c005/src}
\newcommand{\DirFiles}{64}
\newcommand{\DirLines}{5{,}460}
\newcommand{\ConcDecls}{27}
\newcommand{\ConcFiles}{13}
\newcommand{\UnitMembers}{144}
\newcommand{\UnitInDir}{144}
\newcommand{\UnitOutside}{0}
\newcommand{\UnitInNote}{13}
\newcommand{\UnitElsewhere}{131}
\newcommand{\ElsewhereRows}{The Artin map and class groups & 23\\Kummer generators, units and Selmer groups & 30\\Reflection & 8\\Restriction to the cyclotomic kernel & 32\\Definitions of the other subjects & 38\\}
\newcommand{\UnitProofs}{103}
\newcommand{\UnitSteps}{856}
\newcommand{\ConcordanceRows}{\crow{XDual}{\lean{XDual}}{Mk0IdMapCongr.lean}{90--99}
\crow{Dinf}{\lean{Dinf}}{Mk0IdMapCongr.lean}{195--195}
\crow{Xinf}{\lean{Xinf}}{Mk0IdMapCongr.lean}{197--197}
\crow{funRep}{\lean{funRep}}{Mk0IdMapCongr.lean}{121--124}
\crow{conjCond}{\lean{conjCond}}{Mk0IdMapCongr.lean}{103--107}
\crow{EqHom}{\lean{EqHom}}{Mk0IdMapCongr.lean}{109--110}
\crow{Un}{\lean{Un}}{Mk0IdMapCongr.lean}{186--186}
\crow{funRepMem}{\lean{funRepMem}}{FunRepMem.lean}{5--46}
\crow{memPrinc}{\lean{memPrinc}}{DinfNil.lean}{7--12}
\crow{powSucc}{\lean{powSucc}}{DinfNil.lean}{14--29}
\crow{qzTors}{\lean{qzTors}}{DinfNil.lean}{31--43}
\crow{princNhds}{\lean{princNhds}}{DinfNil.lean}{45--60}
\crow{dinfNil}{\lean{dinfNil}}{DinfNil.lean}{62--235}
\crow{G}{\lean{G}}{Statement.lean}{6--6}
\crow{QZ}{\lean{QZ}}{Statement.lean}{15--15}
\crow{cycZ}{\lean{cycZ}}{Global.lean}{90--92}
\crow{Dec}{\lean{Dec}}{Global.lean}{15--15}
\crow{Ine}{\lean{Ine}}{Global.lean}{45--64}
\crow{princ}{\lean{princ}}{LeNuIffPowMemPrinc.lean}{16--16}
\crow{levU}{\lean{levU}}{ConjStableShaFam.lean}{187--187}
\crow{cycU}{\lean{cycU}}{ConjStableShaFam.lean}{189--189}
\crow{conjHom}{\lean{conjHom}}{RestrictionMemInvariantH1.lean}{115--118}
\crow{conjDec}{\lean{conjDec}}{PartsConj.lean}{69--75}
\crow{conjIne}{\lean{conjIne}}{PartsConj.lean}{77--81}
\crow{xUnip}{\lean{xUnip}}{XUnip.lean}{9--27}
\crow{dualUnip}{\lean{dualUnip}}{DualUnip.lean}{9--34}
\crow{shaDGen}{\lean{shaDGen}}{ShaDGen.lean}{9--59}}
\expandafter\newcommand\csname LQ@dinfNil\endcsname{\leanquoteline{0}{theorem UpperTheory.V.Kum.dinfNil :}\leanquoteline{4}{∀ (F Ω : Type) [Field F] [NumberField F] [Field Ω] [Algebra F Ω] [IsAlgClosure F Ω] (l : ℕ) [Fact l.Prime] (N : ℕ) (U : Subgroup (G F Ω)) [U.Normal], IsClosed (U : Set (G F Ω)) → ∀ φ ∈ Dinf l N U, (∃ n, ∀ g ∈ levU l U n, funRep ℤ\_[l] (QZ l) (cycU l U) g φ = φ) ∧}\leanquoteline{4}{∀ γ ∈ levU l U 1, ∃ k : ℕ, (l : ℤ\_[l]) \textasciicircum{} k • φ = 0 ∧ ((funRep ℤ\_[l] (QZ l) (cycU l U) γ - 1) \textasciicircum{} k) φ = 0}}
\expandafter\newcommand\csname LQ@XDual\endcsname{\leanquoteline{0}{def XDual (N : ℕ) (U : Subgroup (G F Ω)) (R M : Type) [CommRing R] [AddCommGroup M] [Module R M] : Submodule R (U → M) where}\leanquoteline{2}{carrier := \{φ | IsLocallyConstant φ ∧ (∀ a b, φ (a * b) = φ a + φ b) ∧}\leanquoteline{4}{(∀ (A : ValuationSubring Ω) (u : U), u.1 ∈ Ine (F := F) A → φ u = 0) ∧}\leanquoteline{4}{∀ A : ValuationSubring Ω, A.valuation (N : Ω) < 1 → ∀ u : U, u.1 ∈ Dec (F := F) A → φ u = 0\}}\leanquoteline{2}{zero\_mem' := ⟨IsLocallyConstant.const \_, by simp, by simp, by simp⟩}\leanquoteline{2}{add\_mem' \{a b\} ha hb := ⟨ha.1.add hb.1, fun x y => by simp only [Pi.add\_apply, ha.2.1, hb.2.1]; abel,}\leanquoteline{4}{fun A u h => by simp [ha.2.2.1 A u h, hb.2.2.1 A u h],}\leanquoteline{4}{fun A hA u h => by simp [ha.2.2.2 A hA u h, hb.2.2.2 A hA u h]⟩}\leanquoteline{2}{smul\_mem' c a ha := ⟨ha.1.comp (c • ·), fun x y => by simp [ha.2.1], fun A u h => by simp [ha.2.2.1 A u h],}\leanquoteline{4}{fun A hA u h => by simp [ha.2.2.2 A hA u h]⟩}}
\expandafter\newcommand\csname LQ@funRepMem\endcsname{\leanquoteline{0}{theorem UpperTheory.V.Kum.funRepMem :}\leanquoteline{4}{∀ (F Ω : Type) [Field F] [NumberField F] [Field Ω] [Algebra F Ω] [IsAlgClosure F Ω] (N : ℕ) (U : Subgroup (G F Ω)) [U.Normal] (R M : Type) [CommRing R] [AddCommGroup M] [Module R M], (∀ g, ∀ φ ∈ XDual N U R M, funRep R M U g φ ∈ XDual N U R M) ∧ (∀ g ∈ U, ∀ φ ∈ XDual N U R M, funRep R M U g φ = φ) ∧}\leanquoteline{4}{∀ (ρ : Representation R (G F Ω) M) (φ : U → M), φ ∈ EqHom N U ρ ↔ φ ∈ XDual N U R M ∧ ∀ g u, ρ g (funRep R M U g φ u) = φ u}}


\linespread{1.08}                       % 11 pt on 14.7 pt
\setlength{\parindent}{0pt}
\setlength{\parskip}{5.5pt plus 1.5pt minus 0.5pt}
\setlength{\abovedisplayskip}{7.5pt plus 2pt minus 3pt}
\setlength{\belowdisplayskip}{7.5pt plus 2pt minus 3pt}
\setlength{\abovedisplayshortskip}{3pt plus 2pt}
\setlength{\belowdisplayshortskip}{6pt plus 2pt minus 2pt}
\emergencystretch=1.5em
\clubpenalty=10000 \widowpenalty=10000 \displaywidowpenalty=10000
\relpenalty=7000 \binoppenalty=9000
\raggedbottom
\pagestyle{fancy}\fancyhf{}\renewcommand{\headrulewidth}{0pt}
\fancyfoot[C]{\sffamily\fontsize{7.6}{9}\selectfont\color{inkfaint}\thepage\,/\,\pageref*{LastPage}}
\setlist{topsep=2pt,partopsep=0pt,parsep=0pt,itemsep=4pt,leftmargin=2.1em,labelsep=0.55em}

\newcommand{\Sha}{\text{\normalfont Ш}}
\newcommand{\Q}{\mathbb{Q}}
\newcommand{\Z}{\mathbb{Z}}
\newcommand{\F}{\mathbb{F}}
\newcommand{\Qbar}{\overline{\Q}_\ell}
\newcommand{\QZ}{\Q_\ell/\Z_\ell}
\newcommand{\len}{\operatorname{len}}
\newcommand{\cok}{\operatorname{cok}}
\newcommand{\End}{\operatorname{End}}
\newcommand{\Hom}{\operatorname{Hom}}
\newcommand{\Gal}{\operatorname{Gal}}
\newcommand{\diag}{\operatorname{diag}}
\newcommand{\kw}[1]{\emph{#1}}
\newcommand{\m}[1]{\mbox{$#1$}}        % a short formula that must not be broken across lines
\newcommand{\flog}{\lfloor\log_\ell Z\rfloor}
\newcommand{\head}[1]{\par\addvspace{17pt plus 4pt minus 2pt}\Needspace{6\baselineskip}%
  \noindent{\fontsize{13.2}{16}\selectfont\bfseries\color{ink}#1}\par\nopagebreak\vspace{-1pt}}
\newcommand{\sect}[2]{\refstepcounter{section}\label{#2}\par\addvspace{19pt plus 4pt minus 2pt}\Needspace{7\baselineskip}%
  \noindent\llap{{\fontsize{13.2}{16}\selectfont\bfseries\color{inkfaint}\thesection}\hspace{9pt}}%
  {\fontsize{13.2}{16}\selectfont\bfseries\color{ink}#1}\par\nopagebreak\vspace{-1pt}}
\newcommand{\subhead}[1]{\par\addvspace{9pt plus 2pt}\Needspace{4\baselineskip}\noindent\textbf{#1}\enspace\ignorespaces}
% where a statement lives in the Lean archive
\newcommand{\leantag}{{\sffamily\bfseries\fontsize{6.3}{8}\selectfont\color{leancol}\addfontfeatures{LetterSpace=7}LEAN}}
\newcommand{\sep}{\unskip\nobreak\hspace{5pt}{\color{inkfaint}\textperiodcentered}\hspace{5pt plus 2pt}\ignorespaces}
\newcommand{\leanline}[1]{\par\nopagebreak\vspace{-2.5pt}{\fontsize{8.9}{12.4}\selectfont\color{inksoft}\raggedright
  \noindent\llap{\leantag\hspace{9pt}}#1\par}\vspace{1pt}}
% a displayed statement in a tinted box: #1 colour, #2 text
\newcommand{\statement}[2]{\par\addvspace{7pt}\begingroup\setlength{\fboxsep}{8.5pt}\setlength{\fboxrule}{0.6pt}%
  \noindent\fcolorbox{#1!30!white}{#1!3!white}{\parbox{\dimexpr\linewidth-2\fboxsep-2\fboxrule\relax}{#2}}%
  \endgroup\par\addvspace{7pt}}
\newcommand{\nt}[2]{\par\noindent\hangindent=1.1em\hangafter=1 #1\enspace{\color{inksoft}#2}\par\vspace{1pt}}
\newcommand{\legdot}[1]{\tikz[baseline=-0.62ex]\fill[#1] (0,0) circle (2.5pt);\hspace{3pt}}
\newcommand{\crow}[4]{\csname CS@#1\endcsname & #2 & \textcolor{inksoft}{#3} & \textcolor{inksoft}{#4}\\}
\newcommand{\leanquoteline}[2]{\par\noindent\hangindent=\dimexpr#1\fontdimen2\font+2.2em\relax\hangafter=1
  \hspace*{#1\fontdimen2\font}#2}
\newcommand{\quoted}[1]{\par\vspace{3pt}{\ttfamily\fontsize{7.5}{10.4}\selectfont\color{leancol}\raggedright
  \setlength{\parskip}{0pt}\leanquote{#1}\par}\vspace{4pt}}
\newcommand{\bx}[3]{{\sffamily\bfseries\fontsize{7.3}{9.2}\selectfont\color{#1!82!black}#2}\\[1.5pt]#3}
\newcommand{\cs}[2]{\expandafter\newcommand\csname CS@#1\endcsname{#2}}

% theorems
\newtheoremstyle{plainx}{7pt plus 2pt}{7pt plus 2pt}{\itshape}{}{\bfseries}{.}{0.5em}{}
\newtheoremstyle{uprightx}{7pt plus 2pt}{7pt plus 2pt}{}{}{\bfseries}{.}{0.5em}{}
\theoremstyle{plainx}
\newtheorem{theorem}{Theorem}[section]
\newtheorem{lemma}[theorem]{Lemma}
\newtheorem{proposition}[theorem]{Proposition}
\newtheorem{corollary}[theorem]{Corollary}
\theoremstyle{uprightx}
\newtheorem{definition}[theorem]{Definition}
\newtheorem{remark}[theorem]{Remark}
\newtheorem{example}[theorem]{Example}
\renewcommand{\qedsymbol}{{\setlength{\fboxsep}{0pt}\setlength{\fboxrule}{0.5pt}\fbox{\rule{0pt}{0.46em}\hspace{0.46em}}}}
\makeatletter
\renewenvironment{proof}[1][Proof]{\par\pushQED{\qed}\normalfont\topsep0pt\partopsep0pt
  \trivlist\item[\hskip\labelsep\itshape #1.]\ignorespaces}{\popQED\endtrivlist\@endpefalse}
\makeatother

% where each declaration of the concordance is treated
\cs{XDual}{Def.~\ref{d:char}} \cs{Dinf}{Def.~\ref{d:char}} \cs{Xinf}{Def.~\ref{d:char}}
\cs{funRep}{§\ref{s:setting}, conjugation} \cs{conjCond}{Prop.~\ref{p:action}(c)} \cs{EqHom}{Prop.~\ref{p:action}(c)}
\cs{Un}{§\ref{s:setting}, the groups} \cs{funRepMem}{Prop.~\ref{p:action}} \cs{memPrinc}{Lemma~\ref{l:four}(a)}
\cs{powSucc}{Lemma~\ref{l:four}(b)} \cs{qzTors}{Lemma~\ref{l:four}(c)} \cs{princNhds}{Lemma~\ref{l:four}(d)}
\cs{dinfNil}{Theorem A} \cs{G}{§\ref{s:setting}} \cs{QZ}{§\ref{s:setting}} \cs{cycZ}{§\ref{s:setting}}
\cs{Dec}{Def.~\ref{d:char}} \cs{Ine}{Def.~\ref{d:char}} \cs{princ}{§\ref{s:setting}, the groups}
\cs{levU}{§\ref{s:setting}, the groups} \cs{cycU}{§\ref{s:setting}, the groups} \cs{conjHom}{§\ref{s:setting}, conjugation}
\cs{conjDec}{Prop.~\ref{p:action}, proof} \cs{conjIne}{Prop.~\ref{p:action}, proof} \cs{xUnip}{§\ref{s:use}}
\cs{dualUnip}{§\ref{s:use}} \cs{shaDGen}{§\ref{s:use}}

\begin{document}
\color{ink}

% ------------------------------------------------------------------------------------- title
{\raggedright
{\sffamily\bfseries\fontsize{7.2}{9}\selectfont\color{leancol}\addfontfeatures{LetterSpace=8}DEFORMALIZATION 2\par}
\vspace{5pt}
{\fontsize{21.5}{25.5}\selectfont\bfseries Local nilpotence over a cyclotomic tower\par}
\vspace{5pt}
{\fontsize{12.4}{16}\selectfont\itshape\color{inksoft} One part of a unit of the Lean archive behind the proof of the
Riemann hypothesis announced by L.\,Alpöge, read back into mathematics\par}}
\vspace{8pt}
{\color{rulegrey}\hrule height 0.6pt}
\vspace{7pt}
\thispagestyle{empty}

\textbf{Summary.} Let $F$ be a number field with algebraic closure~$\Omega$, let $\ell$ be a prime, and let $\kappa$ be
the $\ell$-adic cyclotomic character of \m{G=\Gal(\Omega/F)}. The group $G$ acts by conjugation on the characters of
the kernel of~$\kappa$, with values in~$\QZ$. We show that every locally constant character~$\varphi$ is fixed by
all~$g$ with \m{\kappa(g)\equiv1} modulo~$\ell^n$, for some~$n$, and that for every $\gamma$ with
\m{\kappa(\gamma)\equiv1} modulo~$\ell$ a power of \m{\gamma-1} kills~$\varphi$; the exponent $\ell^nk_0$ suffices,
where $\ell^{k_0}$ kills~$\varphi$. This is the \kw{local nilpotence} that the first note of this
series~\cite{N1} used as an input. By duality it makes $\gamma$ unipotent modulo~$\ell$ on the Iwasawa module.
The proof is a compactness argument in the profinite group~$G$ and one congruence between polynomials.

\statement{leancol}{\fontsize{9.6}{13.2}\selectfont\textbf{What this note is.} On 20 November 2026 L.\,Alpöge
announced a proof of the Riemann hypothesis, with a Lean archive \cite{Al}. That is a claim, and nothing here asserts
that the hypothesis is proved. The reading project divides the part of the archive that the final theorems use
into units. This note treats one part of one unit, on Kummer theory: \UnitInNote{} of its \UnitMembers{}
declarations, those that define the modules of characters and prove local nilpotence. The other \UnitElsewhere{}
belong to notes to come (§\ref{s:notes}). Everything stated here is proved in full. Every statement was compared
with its formal counterpart at commit \href{\PinUrl}{{\leanfont\Pin}} of the public mirror; Appendix~A is the
concordance of \ConcDecls{} declarations. The note was written from the structured proofs of the reading project,
without opening the Lean files of the unit (§\ref{s:notes}).}

% ------------------------------------------------------------------------------------- 1
\sect{The statement}{s:intro}
In the first note~\cite{N1} the lengths of certain groups of Galois cohomology along the cyclotomic tower of a
number field were expressed through one operator on one module, the Iwasawa module~$X_\infty$, and the estimates
needed that this operator is unipotent modulo~$\ell$. The module $X_\infty$ is the dual of a discrete module
$D_\infty$ of characters, and on a discrete module unipotence takes the form proved here: every element is killed
by a power of \m{\gamma-1}. The first note stated this as its Input~1, with a sketch. This note proves it, in the
generality of the archive. The notation is set up in §\ref{s:setting}.

\statement{arith}{\phantomsection\label{t:A}\textbf{Theorem A} (local nilpotence).\enspace\emph{Let $H$ be a closed
normal subgroup of~$G$, let \m{N\ge0} be an integer, and let $\varphi$ be an element of the module
\m{D_\infty=D_N(H_\infty,\QZ)} of characters of \m{H_\infty=H\cap\ker\kappa}.}
\begin{enumerate}[label=\textup{(\roman*)},topsep=3pt,itemsep=1pt]
\item \emph{There is an integer \m{n\ge0} such that \m{g\varphi=\varphi} for every \m{g\in H_n}.}
\item \emph{Some power of $\ell$ kills~$\varphi$. If \m{\ell^{k_0}\varphi=0} and $n$ is as in \textup{(i)}, then for
every \m{\gamma\in H_1} and \m{k=\ell^nk_0}}
\[
(\gamma-1)^{k}\varphi=0\qquad\text{\emph{and}}\qquad\ell^{k}\varphi=0 .
\]
\end{enumerate}}
\leanline{\lean{dinfNil}}

Here $H_n$ is the group of all \m{g\in H} with \m{\kappa(g)\equiv1} modulo~$\ell^n$, and $g\varphi$ is the character
\m{x\mapsto\varphi(g^{-1}xg)}. The archive states (ii) as the existence of one exponent~$k$ with the two
properties; its proof gives the value above.

% ------------------------------------------------------------------------------------- 2
\sect{Characters and conjugation}{s:setting}
\subhead{Conventions.} Throughout, $\ell$ is a prime number, $\Z_\ell$ the ring of $\ell$-adic integers, $\Q_\ell$ its
field of fractions, and $\QZ$ the quotient group, a module over~$\Z_\ell$. Let $F$ be a number field, $\Omega$ an
algebraic closure of~$F$, and \m{G=\Gal(\Omega/F)} the group of automorphisms of $\Omega$ over~$F$, with its usual
topology: the subgroups \m{\Gal(\Omega/E)}, for $E$ a finite extension of $F$ in~$\Omega$, are a basis of
neighbourhoods of~$1$. The group $G$ is profinite: it is compact, and every neighbourhood of~$1$ contains an open
normal subgroup \cite[Ch.\,I, §1]{NSW}. The \kw{$\ell$-adic cyclotomic character} is the continuous homomorphism
\m{\kappa\colon G\to\Z_\ell^\times} to the group of units of~$\Z_\ell$ with \m{g(\zeta)=\zeta^{\kappa(g)}} for every
root of unity $\zeta$ of $\ell$-power order in~$\Omega$.
\leanline{\lean{G}\sep\lean{QZ}\sep\lean{cycZ}}

\subhead{The groups.} For \m{n\ge1} let \m{U_n=1+\ell^n\Z_\ell}, and let \m{U_0=\Z_\ell^\times}; these are open
subgroups of~$\Z_\ell^\times$. For a subgroup $H$ of~$G$ put
\[
H_n=\{g\in H:\ \kappa(g)\in U_n\}\quad(n\ge0),\qquad H_\infty=H\cap\ker\kappa .
\]
So \m{H_0=H}, the groups $H_n$ decrease, and $H_\infty$ is their intersection. If $H$ is normal in~$G$, so are
all of these. For \m{H=G} the group $G_n$ is the Galois group of $\Omega$ over the field generated over~$F$ by
the roots of unity of order~$\ell^n$, and $G_\infty$ that of $\Omega$ over the field $F_\infty$ generated by all
roots of unity of $\ell$-power order.
\leanline{\lean{princ} (the group $U_n$)\sep\lean{levU} (the group $H_n$)\sep\lean{cycU} (the group
$H_\infty$)\sep\lean{Un} (the group $G_n$)}

\begin{definition}\label{d:char}
A \kw{valuation ring} of $\Omega$ is a subring $A$ such that \m{x\in A} or \m{x^{-1}\in A} for every non-zero
\m{x\in\Omega}; its maximal ideal $\mathfrak m_A$ consists of the non-units of~$A$. The \kw{decomposition group}
$G_A$ is the group of all \m{g\in G} with \m{g(A)=A}, and the \kw{inertia group} $I_A$ is the subgroup of those
that act trivially on the residue field~\m{A/\mathfrak m_A}.

Let $C$ be a subgroup of~$G$, let \m{N\ge0} be an integer, and let $M$ be a module over a commutative ring~$R$.
The module \m{D_N(C,M)} consists of the homomorphisms \m{\varphi\colon C\to M} that are locally constant, vanish
on \m{C\cap I_A} for every valuation ring $A$ of~$\Omega$, and vanish on \m{C\cap G_A} for every valuation ring
$A$ with \m{N\in\mathfrak m_A}. It is an $R$-module under pointwise operations. For a subgroup $H$ of~$G$ put
\[
D_\infty=D_N(H_\infty,\QZ),\qquad X_\infty=\Hom_{\Z_\ell}(D_\infty,\QZ).
\]
\end{definition}
\leanline{\lean{Dec}\sep\lean{Ine}\sep\lean{XDual}\sep\lean{Dinf}\sep\lean{Xinf}}

The elements of \m{D_N(C,M)} are the characters of~$C$ that are unramified at every place of~$\Omega$ and trivial
on the decomposition groups of the places above the primes that divide~$N$. For \m{H=G} and \m{N=\ell} the modules
$D_\infty$ and $X_\infty$ are those of the first note \cite[§7]{N1}, with \m{L=F}.

\subhead{Conjugation.} Let $C$ be a normal subgroup of~$G$. For \m{g\in G} and a function \m{\varphi\colon C\to M}
put
\[
(g\varphi)(x)=\varphi(g^{-1}xg)\qquad(x\in C).
\]
This is an $R$-linear action of $G$ on the module of all functions \m{C\to M}: \m{(gh)\varphi=g(h\varphi)}.
\leanline{\lean{conjHom}\sep\lean{funRep}}

\begin{proposition}\label{p:action}
Let $C$ be a normal subgroup of~$G$, let \m{N\ge0}, and let $M$ be a module over a commutative ring~$R$.
\begin{enumerate}[label=\textup{(\alph*)},topsep=2pt,itemsep=2pt]
\item If \m{\varphi\in D_N(C,M)} and \m{g\in G}, then \m{g\varphi\in D_N(C,M)}.
\item If \m{\varphi\in D_N(C,M)} and \m{g\in C}, then \m{g\varphi=\varphi}.
\item Let $\rho$ be a representation of $G$ on~$M$ by $R$-linear maps. An element $\varphi$ of \m{D_N(C,M)}
satisfies \m{\varphi(gxg^{-1})=\rho(g)\varphi(x)} for all \m{g\in G} and \m{x\in C} if and only if
\m{\rho(g)\bigl((g\varphi)(x)\bigr)=\varphi(x)} for all \m{g\in G} and \m{x\in C}.
\end{enumerate}
\end{proposition}
\begin{proof}
(a) The map \m{x\mapsto g^{-1}xg} is a continuous homomorphism of $C$ to itself, so $g\varphi$ is a locally
constant homomorphism. Let $A$ be a valuation ring of~$\Omega$. Then \m{g^{-1}(A)} is a valuation ring with maximal
ideal \m{g^{-1}(\mathfrak m_A)}, and
\[
g^{-1}G_Ag=G_{g^{-1}(A)},\qquad g^{-1}I_Ag=I_{g^{-1}(A)} .
\]
So for \m{x\in C\cap I_A} the element \m{g^{-1}xg} lies in \m{C\cap I_{g^{-1}(A)}}, where $\varphi$ vanishes. If
\m{N\in\mathfrak m_A}, then \m{N=g^{-1}(N)} lies in \m{g^{-1}(\mathfrak m_A)}, and the same argument applies to the
decomposition groups.

(b) For \m{g\in C} the three factors of \m{g^{-1}xg} lie in~$C$, and
\m{\varphi(g^{-1}xg)=-\varphi(g)+\varphi(x)+\varphi(g)=\varphi(x)}, because $M$ is commutative.

(c) Put \m{y=g^{-1}xg}. The second condition reads \m{\rho(g)\varphi(y)=\varphi(gyg^{-1})}, and $y$ runs through $C$
when $x$ does.
\end{proof}
\leanline{\lean{funRepMem}, with \lean{conjDec} and \lean{conjIne}\sep the first condition of (c):
\lean{conjCond}, \lean{EqHom}}

By (a) the group $G$ acts on \m{D_N(C,M)}, and by (b) the action factors through~$G/C$. For \m{C=H_\infty} this is
the action on $D_\infty$ in Theorem~A, and it factors through \m{G/H_\infty}.

% ------------------------------------------------------------------------------------- 3
\sect{Four facts about \texorpdfstring{$\ell$}{l}-adic numbers}{s:four}
\begin{lemma}\label{l:four}
\begin{enumerate}[label=\textup{(\alph*)},topsep=2pt,itemsep=2pt]
\item A unit $u$ of~$\Z_\ell$ lies in the kernel of the reduction of units modulo~$\ell^m$ if and only if $\ell^m$
divides \m{u-1}. So this kernel is~$U_m$.
\item Let \m{v\in\Z_\ell} and \m{j\ge0}. If $\ell^{j+1}$ divides \m{v-1}, then $\ell^{j+2}$ divides \m{v^\ell-1}.
\item Every element of~$\QZ$ is killed by a power of~$\ell$.
\item Every neighbourhood of~$1$ in~$\Z_\ell^\times$ contains $U_n$ for some~$n$.
\end{enumerate}
\end{lemma}
\begin{proof}
(a) The reduction of $u$ is~$1$ exactly when \m{u-1} lies in the kernel $\ell^m\Z_\ell$ of the reduction of~$\Z_\ell$.

(b) Write \m{v-1=\ell^{j+1}c} with \m{c\in\Z_\ell}, and \m{v^\ell-1=(v-1)\,s} with \m{s=1+v+\dots+v^{\ell-1}}. Each
\m{v^i-1} is divisible by \m{v-1}, so \m{s-\ell=\sum_{i<\ell}(v^i-1)=(v-1)\,d} for some \m{d\in\Z_\ell}. Hence
\[
v^\ell-1=\ell^{j+1}c\,\bigl(\ell+\ell^{j+1}c\,d\bigr)=\ell^{j+2}\bigl(c+\ell^{j}c^2d\bigr).
\]

(c) An element of $\QZ$ is the class of some \m{x\in\Q_\ell}, and \m{\ell^mx\in\Z_\ell} for large~$m$.

(d) The topology of $\Z_\ell^\times$ is induced from~$\Z_\ell$, so a neighbourhood of~$1$ contains all units $u$
with \m{u-1\in\ell^n\Z_\ell} for some~$n$.
\end{proof}
\leanline{(a) \lean{memPrinc}\sep (b) \lean{powSucc}\sep (c) \lean{qzTors}\sep (d) \lean{princNhds}}

By (b) and induction, an element $v$ of \m{1+\ell\Z_\ell} satisfies
\begin{equation}\label{e:pow}
v^{\ell^j}\equiv1\pmod{\ell^{j+1}}\qquad\text{for every }j\ge0 .
\end{equation}

% ------------------------------------------------------------------------------------- 4
\sect{Proof of Theorem A}{s:proof}
Write \m{C=H_\infty}. It is a normal subgroup of~$G$, and it is closed, because $H$ is closed and $\kappa$ is
continuous. So $C$ is compact. The proof uses of $\varphi$ only that it is a locally constant homomorphism
\m{C\to\QZ}.

\subhead{Proof of (i).} The kernel of~$\varphi$ is an open subgroup of~$C$, so it is the intersection of $C$ with
an open subset of~$G$ that contains~$1$. This subset contains an open normal subgroup $W$ of~$G$. Thus
\begin{equation}\label{e:w}
\varphi(x)=0\qquad\text{for all }x\in W\cap C .
\end{equation}
\emph{Every \m{c\in C} fixes~$\varphi$.} This is Proposition~\ref{p:action}(b); it uses only that $\varphi$ is a
homomorphism.

\emph{Every \m{w\in W} fixes~$\varphi$.} Let \m{x\in C}. The element \m{w^{-1}xwx^{-1}} lies in~$C$, because $C$ is
normal, and in~$W$, because it is the product of $w^{-1}$ and \m{xwx^{-1}} and $W$ is normal. By \eqref{e:w},
\[
(w\varphi)(x)=\varphi\bigl(w^{-1}xwx^{-1}\cdot x\bigr)=\varphi\bigl(w^{-1}xwx^{-1}\bigr)+\varphi(x)=\varphi(x).
\]
Since \m{(ab)\varphi=a(b\varphi)}, every element of the set
\[
S=\{wc:\ w\in W,\ c\in C\}
\]
fixes~$\varphi$. The set $S$ is a union of cosets of the open subgroup~$W$, so it is open.

Let $K$ be the set of the elements of $H$ that are not in~$S$. It is closed in the compact group~$G$, hence compact,
and its image $\kappa(K)$ is a compact, hence closed, subset of~$\Z_\ell^\times$. It does not contain~$1$: an
element of $H$ with \m{\kappa=1} lies in $C$, and \m{C\subseteq S}. By Lemma~\ref{l:four}(d) there is $n$ such that
$U_n$ does not meet~$\kappa(K)$. Let \m{g\in H_n}. Then \m{\kappa(g)\in U_n}, so \m{g\notin K}; as \m{g\in H}, this
means \m{g\in S}, and $g$ fixes~$\varphi$.

\subhead{Proof of (ii).} The function $\varphi$ is locally constant on a compact space, so it takes finitely many
values. By Lemma~\ref{l:four}(c) each of them is killed by a power of~$\ell$. So \m{\ell^{k_0}\varphi=0} for some
\m{k_0\ge0}.

Let $n$ be as in~(i), put \m{q=\ell^n}, and let \m{\gamma\in H_1}. Then \m{\kappa(\gamma)\in1+\ell\Z_\ell}, and by
\eqref{e:pow} with \m{j=n} the unit \m{\kappa(\gamma^q)=\kappa(\gamma)^{q}} lies in~$U_n$. So \m{\gamma^q\in H_n}, and
by (i)
\begin{equation}\label{e:q}
\gamma^q\varphi=\varphi .
\end{equation}
Let $T$ be the endomorphism \m{\gamma-1} of the $\Z_\ell$-module of all functions \m{C\to\QZ}, and let $J$ be the
set of polynomials \m{P\in\Z[Y]} with \m{P(T)\varphi=0}. It is an ideal of the polynomial ring~$\Z[Y]$. It contains
\[
(Y+1)^q-1\qquad\text{and}\qquad\ell^{k_0},
\]
the first by \eqref{e:q}, since \m{T+1=\gamma}, and the second by the choice of~$k_0$. Now $q$ is a power of~$\ell$,
so \m{(Y+1)^q=Y^q+1} in the polynomial ring over the field with $\ell$ elements. Hence there is \m{Q\in\Z[Y]} with
\[
(Y+1)^q-Y^q-1=\ell\,Q ,\qquad\text{that is,}\qquad Y^q=\bigl((Y+1)^q-1\bigr)-\ell\,Q .
\]
Modulo~$J$ this gives \m{Y^q\equiv-\ell Q}, and then
\[
Y^{qk_0}\equiv(-Q)^{k_0}\,\ell^{k_0}\equiv0\pmod J .
\]
So \m{T^{qk_0}\varphi=0}, which is \m{(\gamma-1)^k\varphi=0} for \m{k=qk_0=\ell^nk_0}. Finally \m{\ell^k\varphi=0}
because \m{k\ge k_0}. \hfill\qedsymbol

\medskip
The exponent does not depend on~$\gamma$. For \m{\ell=2} every unit of $\Z_2$ is congruent to~$1$ modulo~$2$, so
\m{H_1=H} and (ii) holds for every \m{\gamma\in H}.

% ------------------------------------------------------------------------------------- 5
\sect{Where the theorem is used}{s:use}
\emph{Unipotence on the dual.} Let $H$ be a closed normal subgroup of~$G$ and \m{\gamma\in H_1}. By Theorem~A(ii)
for~$\gamma^{-1}$, the operator \m{\gamma^{-1}-1} is locally nilpotent on~$D_\infty$: every element is killed by
a power of it. If $X_\infty$ is finitely generated over~$\Z_\ell$, the proposition on dual modules of the first note
\cite[Prop.\,7.1]{N1} shows that the action of $\gamma$ on~$X_\infty$, \m{(\gamma x)(\varphi)=x(\gamma^{-1}\varphi)},
is unipotent modulo~$\ell$: some power of \m{\gamma-1} maps $X_\infty$ into~$\ell X_\infty$. For \m{H=G} and
\m{N=\ell} this is how Input~1 of the first note enters the proof of its Theorem~D.

\emph{Invariants.} Part (i) is used in the descent to a layer of the tower \cite[§8, Input~2b]{N1}. There one has a
subgroup $V$ of~$G$ and an element \m{\eta\in V} such that $\kappa(V)$ lies in the closure of the set of powers
of~$\kappa(\eta)$, and one needs that an element of $D_\infty$ fixed by~$\eta$ is fixed by~$V$, for the action
twisted by a power of~$\kappa$. By (i) the element is fixed by~$G_n$ for some~$n$. This also holds for the twisted
action when $n$ is at least the exponent~$k_0$ of~(ii), because a unit in~$U_n$ then acts trivially on an element
killed by~$\ell^{k_0}$. And every element of $V$ is the product of a power of~$\eta$ and an element of~$G_n$.
\leanline{\lean{xUnip}, by \lean{dualUnip}\sep\lean{shaDGen}}

% ------------------------------------------------------------------------------------- 6
\sect{Notes on the formal text}{s:notes}
\begin{enumerate}[label=\textup{\arabic*.},leftmargin=1.6em,itemsep=5pt]
\item \emph{The conditions at the places are not used.} The formal statement takes $\varphi$ in~$D_\infty$. Its
proof unfolds this hypothesis once and keeps two of its four clauses: $\varphi$ is locally constant, and it is
additive. The clauses on inertia and decomposition groups are not referred to again, and neither is the
integer~$N$. So Theorem~A holds for every locally constant homomorphism \m{H_\infty\to\QZ}. The conditions matter
for Proposition~\ref{p:action}(a), which makes $D_\infty$ a module with an action of~$G$.
\item \emph{The exponent.} The formal statement asserts that some~$k$ has the two properties of~(ii). The proof
constructs \m{k=\ell^nk_0}, the same for all~$\gamma$; this is the form stated above.
\item \emph{Where the hypotheses on $H$ enter.} Closedness gives the compactness of $H_\infty$, used for the open
subgroup~$W$ and for the finitely many values, and the compactness of the set~$K$. Normality makes $H_\infty$ normal,
so that conjugation acts.
\item \emph{The coefficients.} The module $\QZ$ enters through Lemma~\ref{l:four}(c) alone. The argument of
§\ref{s:proof} applies to locally constant homomorphisms into any module whose elements are killed by powers
of~$\ell$. This remark is ours; the archive states the theorem for~$\QZ$.
\item \emph{The polynomial ideal} is the device of the formal proof, followed here step by step. The sketch in the
first note used the same congruence, applied $k_0$ times to elements; the ideal does this bookkeeping once.
\item \emph{The rest of the unit.} The unit of the archive that holds these declarations has \UnitMembers{}
declarations that the final theorems use. This note treats \UnitInNote. The others are, by source file:

\par\vspace{2pt}
{\fontsize{9.2}{12}\selectfont\noindent\hspace*{1.2em}\begin{tabular}{@{}l@{\hspace{14pt}}r@{}}
\ElsewhereRows
\end{tabular}\par}
\vspace{2pt}

The grouping is provisional. Each group is the subject of a note to come.
\end{enumerate}

\subhead{How this note was written.} The reading project supplied a dossier of the unit: the statements, a
glossary, and a structured proof of every theorem in which each step quotes the lines of the archive that it
renders. A program compared the quotations with the archive, for all \UnitSteps{} steps of the \UnitProofs{}
proofs of the unit, and found that no line of these proofs is left without a step. The proofs of this note were
written from those structured proofs. The Lean files of the unit were not opened. The statement of item~1 rests
on the quoted lines of the structured proof of the theorem: the hypothesis on~$\varphi$ occurs in four of them
and nowhere else. The statements quoted below were read from the archive by the build script, and they agree
with the dossier's.

% ------------------------------------------------------------------------------------- A
\head{Appendix A. Concordance with the Lean archive}
{\fontsize{9.6}{12.8}\selectfont\color{inksoft}\noindent Each name is a link to the lines of the declaration at
commit {\leanfont\Pin} of the public mirror of the archive. The line ranges and the quotations below are read from
the archive by the build script, which also checks them against the SHA-256 of each file. The first thirteen rows
are the declarations of the unit that this note treats.\par}
\vspace{4pt}
{\fontsize{8.6}{11.6}\selectfont\renewcommand{\leanfont}{\ttfamily\fontsize{7.7}{10}\selectfont\color{leancol}}
\setlength{\LTleft}{0pt}\setlength{\LTright}{\fill}\setlength{\tabcolsep}{0pt}
\begin{longtable}{@{}p{30mm}p{56mm}p{43mm}r@{}}
\textbf{In this note} & \textbf{Lean declaration} & \textbf{File} & \textbf{Lines}\\\midrule\endhead
\ConcordanceRows
\end{longtable}}

\Needspace{8\baselineskip}
\subhead{Three statements as they are written in the archive.} Theorem~A; here {\leanfont U} is the group~$H$,
{\leanfont levU l U n} is~$H_n$, {\leanfont cycU l U} is~$H_\infty$, and {\leanfont funRep} is the action by
conjugation:
\quoted{dinfNil}
The module of characters \m{D_N(C,M)}, where {\leanfont U} is the group~$C$:
\quoted{XDual}
Proposition~\ref{p:action}:
\quoted{funRepMem}

% ------------------------------------------------------------------------------------- notation
\head{Notation}
{\fontsize{9.6}{12.6}\selectfont\color{inksoft}\noindent The symbols of the text. A name in brackets is the one
used in the archive.\par}
\vspace{4pt}
\begingroup\fontsize{8.6}{10.9}\selectfont\raggedright
\renewcommand{\leanfont}{\ttfamily\fontsize{7.6}{10}\selectfont\color{leancol}}
\setlength{\columnsep}{9mm}\setlength{\parskip}{0pt}
\begin{multicols}{2}
\nt{$\ell$, $\Z_\ell$, $\Q_\ell$, $\QZ$}{a prime; the $\ell$-adic integers and numbers; their quotient group}
\nt{$F$, $\Omega$, $E$, $F_\infty$}{a number field; an algebraic closure; a finite extension; $F$ with all roots of unity of $\ell$-power order}
\nt{$G$, $G_n$, $G_\infty$}{\m{\Gal(\Omega/F)}; the groups $H_n$, $H_\infty$ for \m{H=G} [{\leanfont Un}]}
\nt{$\kappa$, $\zeta$}{the $\ell$-adic cyclotomic character [{\leanfont cycZ}]; a root of unity}
\nt{$U_n$}{the units \m{1+\ell^n\Z_\ell}, with \m{U_0=\Z_\ell^\times} [{\leanfont princ}]}
\nt{$H$, $H_n$, $H_\infty$}{a subgroup of $G$; its elements with \m{\kappa\in U_n} [{\leanfont levU}]; \m{H\cap\ker\kappa} [{\leanfont cycU}]}
\nt{$C$, $N$, $R$, $M$}{a subgroup of~$G$; an integer \m{\ge0}; a commutative ring; an $R$-module}
\nt{$A$, $\mathfrak m_A$, $G_A$, $I_A$}{a valuation ring of~$\Omega$; its maximal ideal; decomposition and inertia groups}
\nt{$D_N(C,M)$, $D_\infty$, $X_\infty$}{characters of $C$ (Definition~\ref{d:char}) [{\leanfont XDual}]; \m{D_N(H_\infty,\QZ)}; its dual}
\nt{$\varphi$, $g\varphi$}{a character; its conjugate \m{x\mapsto\varphi(g^{-1}xg)}}
\nt{$g$, $h$, $\gamma$, $\eta$, $x$, $y$, $c$, $w$, $a$, $b$}{elements of~$G$ (also: $c$, $d$ elements of $\Z_\ell$ in Lemma~\ref{l:four}; $x$ in $\Q_\ell$)}
\nt{$\rho$}{a representation of $G$ on~$M$}
\nt{$u$, $v$, $s$, $d$}{elements of~$\Z_\ell$; \m{s=1+v+\dots+v^{\ell-1}}}
\nt{$W$, $S$, $K$, $V$}{an open normal subgroup; the set $WC$; the set \m{H\smallsetminus S}; a subgroup of $G$ (§\ref{s:use})}
\nt{$n$, $m$, $j$, $i$, $k$, $k_0$, $q$}{integers; \m{\ell^{k_0}\varphi=0}; \m{q=\ell^n}}
\nt{$T$, $J$, $P$, $Q$, $Y$}{the operator \m{\gamma-1}; the ideal of polynomials that kill~$\varphi$; polynomials; the indeterminate}
\nt{$L$}{the number field of the first note}
\end{multicols}
\endgroup

% ------------------------------------------------------------------------------------- references
\head{References}
\begingroup\fontsize{9.4}{12.6}\selectfont\raggedright\renewcommand{\section}[2]{}
\begin{thebibliography}{NSW}\setlength{\itemsep}{2pt}
\bibitem[Al]{Al} L.\,Alpöge, announcement of proofs of the Riemann hypothesis and of the generalized Riemann
  hypothesis for Dirichlet characters, 20 November 2026, with a Lean archive;
  \href{\PinUrl}{public mirror of the archive at commit {\leanfont\Pin}}.
\bibitem[N1]{N1} \emph{Twisted cokernels in a cyclotomic tower}, Deformalization~1, riemann-rev,
  \href{https://nasqret.github.io/riemann-rev/digests/tower-bound.pdf}{\texttt{digests/tower-bound.pdf}}.
\bibitem[NSW]{NSW} J.\,Neukirch, A.\,Schmidt and K.\,Wingberg, \emph{Cohomology of Number Fields}, 2nd ed.,
  Grundlehren der mathematischen Wissenschaften 323, Springer, 2008.
\end{thebibliography}
\endgroup

\end{document}
